% Notes on this revision, submitted in place of a point-by-point response to reviewers.
% The author does not have the text of the referee reports for Ms. Ref. No.
% JSPI-D-26-00452 at the time of this revision, so this document summarizes the
% revision comprehensively rather than mapping changes to individual numbered comments.
\documentclass[11pt,a4paper]{article}
\usepackage[margin=1in]{geometry}
\usepackage{parskip}
\usepackage{enumitem}
\usepackage{url}
\pagestyle{empty}

\begin{document}

\noindent\textbf{Notes on this revision}\\
Ms.\ Ref.\ No.\ JSPI-D-26-00452 --- ``WAGER: Within-cell Antisymmetric Gain Evaluation
of Resolution''

\vspace{1em}

\noindent This note is not a point-by-point response to the referee reports: at the
time of preparing this revision, the author did not have the text of the individual
reviewer comments in hand. Rather than guess at what was said, this document
summarizes the revision comprehensively --- what was added, what was corrected, and
why --- so the editors and reviewers can see everything that changed regardless of
which points it does or does not happen to address. The author would welcome the
opportunity to prepare a conventional point-by-point response once the reports are
available, if the editors judge that necessary before a decision.

\vspace{1em}
\noindent\textbf{1. New theoretical results}

\begin{enumerate}[leftmargin=1.5em]
\item \textbf{Bregman-score generalization} (Theorem 2, Section 3.3). The covariance
identity previously stated only for the quadratic score is shown to hold for the
entire family of scores generated by a strictly convex, differentiable function on the
simplex, with the quadratic and log-score identities recovered as named corollaries
rather than treated as an unexplained special case and an informal sensitivity check.
\item \textbf{Asymptotic normality} (Theorem 3, Section 3.9). A central limit theorem
for the dataset-level estimator, under a bounded score and a condition ruling out one
cluster from dominating the sample, together with consistency of the already-implemented
sandwich variance estimator. This is the formal justification for reading the
image-clustered intervals reported throughout the paper as asymptotically valid.
\item \textbf{Sensitivity bound for an unrecorded confounder} (Proposition 3 and its
worst-case Corollary, Section 3.4--3.5). A Cauchy--Schwarz argument, in the spirit of
Rosenbaum (2002) and Cinelli \& Hazlett (2020), that bounds how far an unrecorded
covariate could move the reported alignment estimate from what a fully adjusted audit
would find, given a single elicitable correlation parameter, with a data-free
worst-case fallback needing only the label cardinality and score bound.
\item \textbf{Relation to Diebold--Mariano/Giacomini--White} (Corollary, Section 3.9).
At a trivial audit feature (no declared prior structure), the paper's own statistic and
its interval are shown to reduce exactly to a clustered form of the classical
comparative-forecast-evaluation test, situating the method precisely relative to that
literature rather than only by analogy.
\end{enumerate}

\vspace{0.5em}
\noindent\textbf{2. New empirical domain}

A third application, long-tailed text classification on 20 Newsgroups (Section 4.9),
was added alongside the existing scene-graph-generation and long-tailed image
classification studies. This domain has neither a spatial nor a class-frequency prior,
directly testing the estimator's claimed generality on a domain unlike either of the
first two rather than only asserting it. The two domains' training-recipe compositions
turn out to differ --- the same recipe carries the real instance-alignment gain in one
domain and almost none in the other --- and the revision states this as suggestive
pending the fuller (multi-seed, calibration-matched) protocol already applied to the
image-classification domain, not as an already-settled finding.

\vspace{0.5em}
\noindent\textbf{3. Corrections identified during further internal review}

Beyond the additions above, a second internal review pass after the theoretical and
empirical additions were drafted found and corrected the following:

\begin{enumerate}[leftmargin=1.5em]
\item \textbf{A related-work argument was found to be incorrect and was removed.}
Section 2.5 had argued that subtracting two separately-estimated single-model score
resolutions does not recover the paper's own estimate because the relevant finite-sample
bias ``scales with each model's own cell counts.'' On inspection this is false when both
compared models share the same declared audit partition: the same attenuation factor
then applies to each model's own resolution estimate, and subtracting the two recovers
the paper's own estimate of the contrast exactly. The argument was replaced with the
paper's actually-valid distinguishing point --- that the classical literature supplies
no inference for the covariance between two such estimates from the same data, which
this paper's construction targets directly --- and three directly relevant references
were added: Ferro \& Fricker (2012, \emph{QJRMS}), Siegert (2014, \emph{QJRMS}), and
DeLong, DeLong \& Clarke-Pearson (1988, \emph{Biometrics}).
\item \textbf{The new asymptotic normality theorem's proof required an explicit
condition that was missing.} The original proof sketch invoked a within-cell law of
large numbers under a condition that bounds only the aggregate identified sample
fraction, not any individual audit cell's size. An explicit condition was added --- the
audit feature takes values in a fixed, finite set with positive limiting mass per cell,
so every cell's size diverges in the idealized large-sample limit --- and the proof was
corrected to use it. The manuscript now states plainly that this is an asymptotic
idealization the real scene-graph application has not attained in-sample (most audit
cells there hold only the minimum of two examples), and that the reported empirical
interval-coverage simulation, not the theorem directly, is the finite-sample evidence
for the interval's practical validity.
\item \textbf{The headline empirical claim was restated to reflect a robustness check
the paper itself runs.} The audit of a released scene-graph debiasing method (Section
4.6) reports the same comparison under two proper scores. The quadratic score's
calibration-matched interval crosses zero; the log score's does not (a fact already
computed but not previously carried into the manuscript's stated conclusion or its own
table). The revision now states precisely what the two scores agree on --- that the
method's improvement is overwhelmingly attributable to fitting the prior rather than to
individual-instance evidence --- and what they disagree on --- whether any
instance-alignment remainder is exactly zero --- rather than reporting only the reading
that happened to show a null result.
\end{enumerate}

\vspace{0.5em}
\noindent\textbf{4. Reproducibility}

All code, cached model outputs, and scripts needed to reproduce every number, table, and
figure in the revised manuscript are available at
\url{https://github.com/vinhqdang/WAGER}, including a script that checks every reported
value against the committed results files.

\end{document}
